\begin{center}
	\zihao{4}\heiti
	% 中国科学院大学\\
	\LARGE{$\boldsymbol{2007}$年硕士学位研究生入学统一考试试题\\
	电动力学B卷}

\end{center}
\noindent
考生须知：\\
1.本试卷满分为150分,全部考试时间总计180分钟。\\
2.所有答案必须写在答题纸上,写在试题纸上或草稿纸上一律无效。

\vspace*{2ex}
\hrule
\vspace*{4ex}

\fancypagestyle{mypagestyle}{%
	\fancyhf{}% Clear header/footer
	\fancyfoot[OR]{第~\thepage~页，共~\pageref{LastPage}~页}%
	\fancyfoot[EL]{第~\thepage~页，共~\pageref{LastPage}~页}%
	\fancyfoot[OL]{科目名称：电动力学 }%
	\fancyfoot[ER]{科目名称：电动力学 }%
	\renewcommand{\headrulewidth}{0pt}% Header rule of .4pt
	\renewcommand{\footrulewidth}{0.5pt}% Header rule of .4pt
}
\fancypagestyle{plain}{%
	\fancyhf{}% Clear header/footer
	\fancyfoot[OR]{第~\thepage~页，共~\pageref{LastPage}~页}%
	\fancyfoot[EL]{第~\thepage~页，共~\pageref{LastPage}~页}%
	\fancyfoot[OL]{科目名称：电动力学 }%
	\fancyfoot[ER]{科目名称：电动力学 }%
}

\setcounter{page}{1}%设置页面计数器
\pagestyle{mypagestyle}

\onehalfspacing
% 排版结束，下面是题目
\vspace{0.4cm}
\noindent
一、（共25分）问答题

\noindent
1）\hspace{0.2em}在相对论电动力学中，试写出由电磁场量构成的两个洛伦兹不变的量

\noindent
2）\hspace{0.2em}写出电偶极辐射的场分布公式

\noindent
3）\hspace{0.2em}写出电磁场矢势和标势满足的适用于一般规范的方程组。\quad 若采用洛伦兹规范，此方程组变为什么形式？

\vspace{0.5cm}
\noindent
二、（共35分）选择题

\noindent
1）\hspace{0.2em}有一内外半径分别为a，b的导体球壳，在其中心处有一点电荷q，在静电平衡状态下在外表面有电荷密度为$\sigma$的均匀电荷分布，则在$r > b$处的电场为:

\vspace{0.2cm}
(a)\hspace{0.6em} 0 \hspace{1.4em}
(b)\hspace{0.6em}$ \displaystyle{\frac{\sigma b^2}{\boldsymbol{\varepsilon_0} r^2} \Vec{e_r}}$\hspace{1.4em}
(c)\hspace{0.6em} $\displaystyle{\frac{q}{\boldsymbol{\varepsilon_0} r^2} \Vec{e_r}}$\hspace{1.4em}
(d)\hspace{0.6em} $\displaystyle{\frac{q}{4 \pi \boldsymbol{\varepsilon_0} r^2} \Vec{e_r}}$
\vspace{0.2cm}

\noindent
2）\hspace{0.2em}两个相等的电荷+$q$相距2d，另有一个接地导体球放在它们中间。如果这两个电荷所受的总静电力等于0，试问球的半径约为多大？

\vspace{0.2cm}
(a)\hspace{0.6em} d/2 \hspace{1.8em}
(b)\hspace{0.6em} d/4 \hspace{1.8em}
(c)\hspace{0.6em} d/6 \hspace{1.8em}
(d)\hspace{0.6em} d/8 \hspace{1.8em}

% \begin{wrapfigure}{r}[0cm]{0pt}
% \includegraphics[width=5.2cm]{test/2007/pic_4.png} 
% \end{wrapfigure}

\vspace{0.2cm}

\noindent
3）\hspace{0.2em}如图所示是一无限长中空圆柱形导体壳的横截面，外半径为b，内半径为c，中间是一横截面半径为a的圆柱形导体，三个圆柱面同轴。此圆柱导体和壳之间为真空。若在导体轴和中心圆柱体上都有相同轴向的均匀电流流动，电流密度都为$\Vec{j}$，导体的磁导率都为$\mu$，则图中A点（距轴线的距离为$\mathrm{r_1}$）和球壳内B点（距轴线的距离为$\mathrm{r_2}$）的磁感应强度大小分别为

\vspace{0.5cm}
(a)\hspace{0.6em} $B_A = 0$, \quad $B_B = \displaystyle{\frac{\mu [ja^2 + j{r_2}^2 - jc^2]}{2r_2} }$
\vspace{0.5cm}

(b)\hspace{0.6em} $B_A = \displaystyle{\frac{j \mu_0 a^2}{2r_1}}$, \quad $B_B = \displaystyle{\frac{\mu [ja^2 + j{r_2}^2 - jc^2]}{2r_2} }$
\vspace{0.5cm}

(c)\hspace{0.6em} $B_A = \displaystyle{\frac{j \mu_0 a^2}{2r_1}}$, \quad $B_B = \displaystyle{\frac{\mu [j{r_2}^2 - jc^2]}{2r_2} }$ 
\vspace{0.5cm}

(b)\hspace{0.6em} $B_A = 0$, \quad  $B_B = \displaystyle{\frac{\mu [ja^2 + jb^2 - jc^2]}{2r_2} }$ 

\vspace{-6cm}

\begin{figure}[h!]
    \flushright
    \includegraphics[height=5.5cm]{test/2007/pic_4.png}
\end{figure}
\vspace{0.2cm}

\noindent
4）\hspace{0.2em}边长为1cm的正立方形金属壁空腔，若该谐振腔中的本征频率在$1.5 \times 10^{15} - 1.6 \times 10^{15}$Hz之间，那么可能存在的电磁驻波模式的数目约为：

\vspace{0.2cm}
(a)\hspace{0.6em} $10^{13}$\hspace{1.4em}
(b)\hspace{0.6em} $10^{14}$\hspace{1.4em}
(c)\hspace{0.6em} $10^{15}$\hspace{1.4em}
(d)\hspace{0.6em} $10^{16}$

\vspace{0.2cm}

\noindent
5）\hspace{0.2em} 麦克斯韦电磁理论适用的范围为

\vspace{0.2cm}
(a)\hspace{0.6em} 宏观不限速\hspace{1.4em}
(b)\hspace{0.6em} 宏观高速\hspace{1.4em}

\vspace{0.2cm}

(c)\hspace{0.6em} 宏观和微观\hspace{1.4em}
(d)\hspace{0.6em} 宏观低速

\newpage
\noindent
三、\hspace{0.2em}（共30分）如图所示有一小型空心带电圆盘，内半径为$R_1$，外半径为$R_2$，盘的厚度可以忽略不计，盘上带有均匀的电荷，面电荷密度为$\sigma$，如果这个圆盘以角速度$\boldsymbol{\omega} = C \sin \alpha t$绕过其中心的垂直轴线转动，$\alpha$为一常数。\vspace{0.2cm}

\noindent1)\hspace{0.6em}试求中心垂直轴线上任一点A的磁感应强度，其中A点距离圆盘中心O点的距离为x

\noindent2)\hspace{0.6em}如果在A点上平行于圆盘放一半径非常小的线圈，半径为b，线圈内任意点的磁感应强度可以认为等于A点的磁场在某一时刻t将伏特计接到线圈上，问伏特计的瞬时读数是多少?

\vspace{-0.4cm}
\begin{figure}[h!]
    \flushright
    \includegraphics[height=6.5cm]{test/2007/pic_5.png}
\end{figure}

\vspace{-0.4cm}
\noindent
四、\hspace{0.2em}（共30分）
有一带电荷q的粒子，在xy平面绕半径为a的圆做圆周运动，角频率为$\boldsymbol{ \omega }$，且满足$\boldsymbol{\omega} a \ll c$，求1)辐射场在球坐标系中的分布，2)平均辐射功率，3)分析其在z方向和xy平面内各方向的偏振特性\vspace{0.2cm}

\hspace{0.4em}(提示：直角坐标和球坐标单位矢量换算如下

\vspace{0.2cm}
\hspace{0.4em}$\displaystyle{
\Vec{e_x} = \sin \theta \cos \phi  \, \Vec{e_r} + \cos \theta \cos \phi \, \Vec{e_{\theta}} - \sin \phi \,  \Vec{e_{\phi}}
}$

\vspace{0.2cm}
\hspace{0.4em}$\displaystyle{
\Vec{e_y} = \sin \theta \sin \phi \,  \Vec{e_r} + \cos \theta \sin \phi \, \Vec{e_{\theta}} + \cos  \phi \,  \Vec{e_{\phi}}
}$

\vspace{0.2cm}
\hspace{0.4em}$\displaystyle{
\Vec{e_z} = \cos \theta \,  \Vec{e_r} - \sin \theta \, \Vec{e_{\theta}}
}$)

\vspace{4cm}

\noindent
五、\hspace{0.2em}（共30分）有一无穷长直导线载有电流I，另有一段直导线AB与电流I共面，并与I垂直，A，B两端到I的距离分别为a和b(b>a)，这段导线以匀速v平行于电流I运动。\vspace{0.2cm}

\noindent1)\hspace{0.6em}试求导线A，B两端的电势差

\noindent2)\hspace{0.6em}在跟随导线AB运动的坐标系里，试求电流I所产生的电场强度和导线A，B两端的电势差

